% 已查地图（先查后写）：查 C-230（最终闭合）、C-229（勘误＋覆盖修补）、C-227（B5）、C-226（D-B）、C-224（T13-A C 严格模板）。
% D0: 本档对象 = 阻尼 M=3 常数的英文正式定理/证明版本（Theorem A + Lemma 1--5 + 复现 + caveat）—— 关系 = 正式化与冻结登记
% D1: 0
% FREEZE-ACK: 本档即冻结期内的正式化与登记（依 §8.1；不产候选结论）

\documentclass[11pt]{article}
\usepackage[margin=1in]{geometry}
\usepackage{amsmath,amssymb,amsthm}
\usepackage{booktabs}
\usepackage[hidelinks]{hyperref}

\newtheorem{theorem}{Theorem}
\newtheorem{lemma}[theorem]{Lemma}
\newtheorem{proposition}[theorem]{Proposition}
\newtheorem{corollary}[theorem]{Corollary}
\theoremstyle{remark}
\newtheorem{remark}[theorem]{Remark}
\newcommand{\CA}{\textsf{[CA]}}
\newcommand{\R}{\mathbb{R}}
\newcommand{\Dd}{\mathbb{D}}
\newcommand{\Om}{\Omega}
\newcommand{\hd}{\operatorname{hd}}
\newcommand{\far}{\operatorname{far}}

\title{The damped three-point cosine min--max constant:\\ a certified identification}
\author{Hui Tang}
\date{Draft v1 --- \today}

\begin{document}
\maketitle

\begin{abstract}
For $z_j=r_j e^{i\varphi_j}$ with $|z_j|\le 1$ and $\max_j|z_j|=1$, consider the
\emph{short-window} real-part power sum $\max_{1\le k\le 5M}\operatorname{Re}\sum_{j\le M}z_j^k$.
For $M=3$ the worst configuration is \emph{damped}, i.e.\ it lies in the interior of the polydisc.
We determine the corresponding constant exactly:
\[
C_3 \;=\; 0.373\,091\,892\,895\,816\,4\ldots
\]
and show that, in the domain of definition $\Om=[0,1]^2\times[0,\pi]^3$, the minimiser set consists of
exactly two points, $z_*$ and its partner $\sigma z_*$ under the paired swap of the two damped factors.
The proof is a three-piece splice: a certified interval Krawczyk enclosure of a six-branch nonsmooth KKT
root (existence, uniqueness, strict positivity of all multipliers, and exact tangency at the root); a
strict local-growth lemma centred at that root; and a certified global branch-and-bound exclusion on the
complement of two balls. All computer-assisted statements carry an explicit interval certificate and a
structurally independent reproduction.
\end{abstract}

\tableofcontents

\section{Setting and statement of the theorem}

\subsection{Definitions}

Throughout, $M=3$, $K:=5M=15$, and the \emph{domain} is the five-dimensional box
\begin{equation}
\Om \;:=\; [0,1]^{2}\times[0,\pi]^{3}\;\subset\;\R^{5},
\qquad x=(r_2,r_3,\varphi_1,\varphi_2,\varphi_3).
\end{equation}
For $1\le k\le K$ define
\begin{equation}
S_k(x) \;:=\; \cos(k\varphi_1)+r_2^k\cos(k\varphi_2)+r_3^k\cos(k\varphi_3),
\qquad
F(x) \;:=\; \max_{1\le k\le K} S_k(x),
\qquad
C_3 \;:=\; \inf_{x\in\Om}F(x).
\end{equation}
The \emph{active set} at $x$ is $A(x):=\{k\le K: S_k(x)=F(x)\}$, and the \emph{active spread} is
\begin{equation}\label{eq:spread}
\delta_A(x) \;:=\; F(x)-\min_{k\in A(x)}S_k(x)\;\ge\;0 .
\end{equation}
The \emph{paired swap} is the involution
\begin{equation}\label{eq:swap}
\sigma:(r_2,r_3,\varphi_1,\varphi_2,\varphi_3)\longmapsto(r_3,r_2,\varphi_1,\varphi_3,\varphi_2),
\qquad C_2:=\{\mathrm{id},\sigma\} .
\end{equation}

\paragraph{Provenance of the parametrisation.}
Writing $z_j=r_je^{i\varphi_j}$ one has $\operatorname{Re}z_j^k=r_j^k\cos(k\varphi_j)$. Normalising the
largest modulus to $1$ and placing it first yields $r_1=1$, so the three-point, linear-window problem
$\max_{1\le k\le 5M}\operatorname{Re}\sum_{j\le M}z_j^k$ reduces to the study of $F$ on $\Om$. Since
$\cos$ is even and $2\pi$-periodic, the angle reductions $\varphi_j\mapsto-\varphi_j$ and
$\varphi_j\mapsto 2\pi-\varphi_j$ make $[0,\pi]$ a fundamental domain in each angle coordinate; no
generality is lost.

\paragraph{The certified box.}
$X_0\subset\R^{11}$ denotes the interval Krawczyk box produced in Lemma~\ref{lem:kkt} (seed radius
$r_0=10^{-12}$, final coordinate width $\le 10^{-21}$), $\hd(X_0)$ its half-diagonal, and $z_*$ the unique
point of $\Om$ solving the KKT system of Lemma~\ref{lem:kkt} inside $X_0$. We write $c_0$ for the decimal
centre of $X_0$, used as the numerical centre of the exclusion balls in Lemma~\ref{lem:global}. Finally,
$B(c,\rho)$ denotes the closed Euclidean ball in $\R^5$, $\far(X,c):=\max_{y\in X}\|y-c\|$ (attained at a
vertex of $X$), and $\underline F_{\mathrm{IA}}(X)$ an interval-arithmetic lower bound for $\min_X F$.

\subsection{Main theorem}

\begin{theorem}[Global minimisation for damped $M=3$]\label{thm:A}
With $\Om,F,C_3$ as above,
\[
\boxed{\;C_3 \;=\; F(z_*) \;=\; 0.373\,091\,892\,895\,816\,4\ldots\;}
\]
and the minimiser set in the normalised domain is
\[
\boxed{\;\operatorname*{arg\,min}_{x\in\Om}F(x)\;=\;\{z_*,\;\sigma z_*\}\;},
\]
two distinct points of $\Om$.
\end{theorem}

The proof is given in \S\ref{sec:proof}; it is the splice of Lemma~\ref{lem:kkt}
(the certified KKT root), Lemma~\ref{lem:sep} (active/non-active separation),
Lemma~\ref{lem:cover} (a uniform first-order covering estimate), Lemma~\ref{lem:growth}
(strict local growth at the root) and Lemma~\ref{lem:global} (certified global exclusion).

\begin{remark}[Minimiser set: normalised domain versus raw permutations]\label{rem:frozen}
The minimiser set \emph{in the normalised domain $\Om$} consists of the two points $z_*$ and $\sigma z_*$.
The corresponding \emph{raw} permutation orbit contains six configurations, whose images under the
normalisation map $\mathcal N$ reduce to these two points, i.e.
$|\mathcal O_{\mathrm{raw}}|=6$ while $|\mathcal N(\mathcal O_{\mathrm{raw}})|=2$
(Proposition~\ref{prop:orbit}). Statements about the number of minimisers always refer to $\Om$.
\end{remark}

\section{Lemma 1: certified KKT root and active-set positivity \CA}\label{lem:kkt}

Let $A_*:=\{1,2,3,4,5,15\}$ and consider the system in the eleven unknowns
\[
z=(r_2,r_3,\varphi_1,\varphi_2,\varphi_3,\lambda_1,\lambda_2,\lambda_3,\lambda_4,\lambda_5,\lambda_{15}),
\]
\begin{equation}\label{eq:kkt}
G(z)\;=\;
\begin{pmatrix}
\sum_{k\in A_*}\lambda_k\,\nabla S_k \\[2pt]
S_1-S_2 \\ S_1-S_3 \\ S_1-S_4 \\ S_1-S_5 \\ S_1-S_{15} \\[2pt]
\sum_{k\in A_*}\lambda_k-1
\end{pmatrix}
\;=\;0
\qquad(5+5+1=11\ \text{equations}).
\end{equation}

\begin{lemma}\label{lem:kkt}
There is exactly one solution $z_*$ of \eqref{eq:kkt} with coordinate vector in the box $X_0$.
At that solution
\[
\lambda_{\min}=\min_{k\in A_*}\lambda_k = 0.111\,860\,12 >0,
\qquad
\|I-YJ(X_0)\|_\infty = 1.965\,1\times10^{-10} < 1,
\]
and every coordinate width of $X_0$ is $\le 10^{-21}$. Consequently the exact root satisfies the six
tangency equations, i.e.\ $S_k(z_*)=F(z_*)$ for all $k\in A_*$, and hence
\begin{equation}\label{eq:deltaA0}
\delta_A(z_*)\;=\;0 .
\end{equation}
\end{lemma}

\begin{proof}[Certification]
The system \eqref{eq:kkt} is solved by an interval Krawczyk operator
\[
K(X)\;=\;m-Y\,G(m)+\bigl(I-Y\,J(X)\bigr)(X-m),\qquad m=\operatorname{mid}(X),
\]
where $J$ is the analytic Jacobian and $Y$ an approximate inverse of $J$ at the high-precision root. Two
\emph{separate} operator conditions are verified in interval arithmetic:
\[
K(X)\subset\operatorname{int}X \quad\text{(existence of a solution in $X$);}\qquad
\|I-YJ(X)\|_\infty<1 \quad\text{(uniqueness).}
\]
Existence and uniqueness are stated separately for this reason. The seed radius is $r_0=10^{-12}$; the
system is ill-conditioned ($\operatorname{cond}\approx2\times10^4$, multipliers spanning $0.112$--$0.816$,
frequencies reaching $15$) and with $r_0=10^{-6}$ one gets $\|I-YJ\|\approx6>1$, so the contraction test
fails --- the small seed is necessary, not cosmetic. The six multipliers certified in $X_0$ are
\[
(\lambda_1,\lambda_2,\lambda_3,\lambda_4,\lambda_5,\lambda_{15})
=(0.200\,054,\ 0.178\,488,\ 0.165\,510,\ 0.194\,819,\ 0.111\,860,\ 0.149\,270),
\]
all strictly positive. Since the solution of \eqref{eq:kkt} inside $X_0$ is unique and the interval
enclosures of the five tie differences $S_1-S_k$ ($k=2,3,4,5,15$) all contain $0$, the unique root
satisfies the ties exactly; together with $\sum_k\lambda_k=1$ and $\lambda_k>0$ this gives $S_k(z_*)=F(z_*)$
for every $k\in A_*$, whence \eqref{eq:deltaA0}.
\end{proof}

\begin{remark}
Containment of $0$ in the interval enclosure of a tie \emph{over the whole box} does not assert that the
ties hold identically on $X_0$. Only the equality at the unique root, \eqref{eq:deltaA0}, is used below.
\end{remark}

\section{Lemma 2: active/non-active separation \CA}\label{lem:sep}

\begin{lemma}\label{lem:sep2}
On the reference box $X_{\mathrm{ref}}$ (the variable part of $X_0$),
\begin{equation}\label{eq:Deltaref}
\Delta_{\mathrm{ref}}
\;:=\;\min_{k\in A_*}\inf_{X_{\mathrm{ref}}}S_k\;-\;\max_{k\notin A_*}\sup_{X_{\mathrm{ref}}}S_k
\;=\;0.279\,596\,813\,6 \;>\;0 .
\end{equation}
With the Lipschitz constants of the two families on $\Om$,
\[
L_{\mathrm{act}}:=\max_{k\in A_*}k\sqrt5=15\sqrt5=33.541\,02,
\qquad
L_{\mathrm{non}}:=\max_{k\notin A_*}k\sqrt5=14\sqrt5=31.304\,952,
\]
the active set is constant on the ball $B\bigl(c_0,\rho_{\mathrm{iso}}\bigr)$ with
\begin{equation}\label{eq:rhoiso}
\rho_{\mathrm{iso}}\;=\;\frac{\Delta_{\mathrm{ref}}}{L_{\mathrm{act}}+L_{\mathrm{non}}}
\;=\;4.311\,71\times10^{-3}.
\end{equation}
\end{lemma}

\begin{proof}[Certification]
For a box $X$, $\inf_X S_k$ and $\sup_X S_k$ are evaluated \emph{exactly per coordinate}: the extrema of
$\cos(k\varphi)$ on an interval are attained at an endpoint or at an odd (resp.\ even) multiple of $\pi$
inside it, and the damping factor $r^k$ enters as an \emph{interval product}. The Lipschitz bound
$\|\nabla S_k\|\le k\sqrt5$ holds because each of the five partial derivatives is bounded by $k$
(displayed: $\partial_{r_j}(r_j^k\cos k\varphi_j)=kr_j^{k-1}\cos k\varphi_j$ and
$\partial_{\varphi_j}(r_j^k\cos k\varphi_j)=-kr_j^k\sin k\varphi_j$, all $\le k$ in modulus). Separation
\eqref{eq:Deltaref} then forces the active set to be constant on any ball of radius at most
$\Delta_{\mathrm{ref}}/(L_{\mathrm{act}}+L_{\mathrm{non}})$, giving \eqref{eq:rhoiso}.
\end{proof}

\begin{remark}[Why the interval product is indispensable]
Bounding $r^k\cos(k\varphi)$ as ``a positive factor times a cosine'' by treating the two factors
independently gives the spurious value $\Delta_{\mathrm{ref}}=-0.522<0$, i.e.\ a false failure. The
genuine product bound yields \eqref{eq:Deltaref}.
\end{remark}

\section{Lemma 3: uniform first-order covering estimate \CA}\label{lem:cover}

\begin{lemma}\label{lem:cover2}
At the certified root $z_*$ the origin lies in the interior of the convex hull of the active gradients,
and, after accounting for the data ball $B(c_0,\rho_{\mathrm{data}})$ with $\rho_{\mathrm{data}}=2.05\times10^{-3}$,
the covering constant
\begin{equation}\label{eq:cX}
c_X \;=\; c-\varepsilon_{\mathrm{pert}}
\;=\;0.302\,091\,5-0.203\,2\;=\;0.098\,812\,009\,77\;>\;0
\end{equation}
satisfies
\[
\max_{k\in A_*}\langle\nabla S_k(x),u\rangle\;\ge\;c_X\qquad\text{for all }x\in B(c_0,\rho_{\mathrm{data}})
\text{ and all }\|u\|=1 .
\]
\end{lemma}

\begin{proof}[Certification]
The six points $\{\nabla S_k(z_*)\}_{k\in A_*}\subset\R^5$ span a simplex. The quantity
$c:=\min_{\|u\|=1}\max_{k\in A_*}\langle\nabla S_k(z_*),u\rangle$ equals the largest ball about the origin
contained in that simplex, i.e.\ the minimum, over its six facets, of the distance from the origin to the
facet hyperplane. The facet normals are obtained from the null space of the corresponding $5\times5$ facet
matrix, and a containment certificate (all remaining vertices strictly on one side of each facet) verifies
that the inscribed ball is indeed interior. Propagating the enclosing data ball of radius
$\rho_{\mathrm{data}}$ by the maximal gradient variation $\varepsilon_{\mathrm{pert}}=0.203\,2$ yields the
usable constant $c_X$ in \eqref{eq:cX}.
\end{proof}

\section{Lemma 4: strict local growth at the certified root \CA}\label{lem:growth}

\begin{lemma}\label{lem:growth2}
Let
\begin{equation}\label{eq:R}
R\;:=\;\max_{k\in A_*}\ \sup_{x\in B(c_0,\rho_{\mathrm{data}})}\bigl\|\nabla^2S_k(x)\bigr\|_2
\;=\;99.899\,695\,515\,930\,13,
\qquad
\rho_g\;:=\;\frac{2c_X}{R}\;=\;1.978\,224\,4\times10^{-3}.
\end{equation}
Then for every $\delta$ with $0<\|\delta\|\le\rho_g$,
\begin{equation}\label{eq:locgrowth}
\boxed{\;F(z_*+\delta)\;\ge\;F(z_*)+c_X\|\delta\|-\frac{R}{2}\|\delta\|^{2}\;>\;F(z_*)\;}
\end{equation}
\end{lemma}

\begin{proof}
Fix $x\in B(c_0,\rho_{\mathrm{data}})$. By Taylor's theorem and \eqref{eq:R},
\[
S_k(x+\delta)\;\ge\;S_k(x)+\langle\nabla S_k(x),\delta\rangle-\frac{R}{2}\|\delta\|^2
\qquad(k\in A_*).
\]
Taking the maximum over $k\in A_*$ and using the elementary fact
\begin{equation}\label{eq:subadditive}
\max_{k}\bigl(a_k+b_k\bigr)\;\ge\;\min_{k}a_k+\max_{k}b_k
\end{equation}
(valid for arbitrary finite families; note that $\max$ is \emph{sub}additive, so the reverse inequality is
false) we obtain the \emph{spread form}
\begin{equation}\label{eq:spreadform}
\boxed{\;F(x+\delta)\;\ge\;F(x)-\delta_A(x)+c_X\|\delta\|-\frac{R}{2}\|\delta\|^{2}\;}
\end{equation}
where $\delta_A(x)$ is the active spread \eqref{eq:spread}. Centred at the root $x=z_*$, the identity
\eqref{eq:deltaA0} of Lemma~\ref{lem:kkt} gives $\delta_A(z_*)=0$, and \eqref{eq:spreadform} becomes
\eqref{eq:locgrowth}. Finally $\rho_g=2c_X/R$ is precisely the (smaller) root of
$c_X\rho-\tfrac R2\rho^2=0$, so the right-hand side of \eqref{eq:locgrowth} is strictly positive for all
$0<\|\delta\|<\rho_g$ and equals $F(z_*)$ only in the limit $\|\delta\|\to\rho_g$.
\end{proof}

\begin{remark}[The spread term is the true logical consumer of the tangency]
If \eqref{eq:spreadform} were available with $\min_{k\in A_*}S_k(x)$ at an \emph{arbitrary} point $x$ of
the reference box, then every such $x$ would be a strict local minimiser of $F$, which is impossible. The
honest statement therefore carries the spread $\delta_A(x)$, and Lemma~\ref{lem:kkt} is consumed at exactly
one place: the identity $\delta_A(z_*)=0$.
\end{remark}

\begin{remark}[Self-consistency of the reference ball]\label{rem:selfcons}
The constants $c_X$ and $R$ must be valid on the whole ball on which \eqref{eq:locgrowth} is asserted.
This is a \emph{proof condition}, not an implementation detail:
\begin{equation}\label{eq:selfcons}
\boxed{\;\rho_{\mathrm{data}}\;\ge\;\rho_g+\hd(X_0)\;}
\end{equation}
Indeed $c_X$ is obtained by propagating the data ball $B(c_0,\rho_{\mathrm{data}})$, while
\eqref{eq:locgrowth} is applied on $B(z_*,\rho_g)$; since $\|z_*-c_0\|\le\hd(X_0)$, \eqref{eq:selfcons} is
exactly what makes the application legitimate. Here
$\rho_{\mathrm{data}}=2.05\times10^{-3}\ge1.978\,224\,4\times10^{-3}+10^{-21}$. A configuration with
$\rho_{\mathrm{data}}=10^{-5}$ and certified radius $4.3\times10^{-3}$ violates \eqref{eq:selfcons} and was
discarded.
\end{remark}

\section{Lemma 5: certified global exclusion \CA}\label{lem:global}

\begin{lemma}\label{lem:global2}
Let $\rho_g$ be as in \eqref{eq:R}, put $B_1:=B(z_*,\rho_g)$, $B_2:=B(\sigma z_*,\rho_g)$, and
\begin{equation}\label{eq:TC}
T_C\;:=\;\sup_{X_0}F+10^{-9}\;=\;F(z_*)+10^{-9}\;=\;0.373\,091\,893\,895\,817 .
\end{equation}
Then $F(x)\ge T_C$ for every $x\in\Om\setminus(B_1\cup B_2)$.
\end{lemma}

\begin{proof}[Certification]
Run an adaptive branch-and-bound on $\Om$ with the exact separable box lower bound
\[
\mathrm{LB}(X)=\max_{1\le k\le K}\Bigl[\min_{I_1}\cos k\varphi_1+\min_{I_2}\bigl(r_2^k\cos k\varphi_2\bigr)
+\min_{I_3}\bigl(r_3^k\cos k\varphi_3\bigr)\Bigr]\;\le\;\min_X F ,
\]
the damped terms again evaluated as interval products. A box $X$ is
\emph{discarded} if $\far(X,c_0)\le\rho_g-10^{-9}$ from one of the two centres --- then
$X\subseteq B(c_0,\rho_g-10^{-9})\subseteq B(z_*,\rho_g)$ because $\hd(X_0)\le10^{-21}$, so $X$ lies in the
region certified by Lemma~\ref{lem:growth2}; \emph{certified} if
$\underline F_{\mathrm{IA}}(X)\ge T_C$; and \emph{split} (bisection along the widest coordinate) otherwise,
unless the box is already below the minimum width $10^{-7}$, in which case it is counted
\emph{unresolved} and would invalidate the lemma. The outcome is:
\begin{center}
\begin{tabular}{lrl}
\toprule
boxes evaluated & $N_{\mathrm{eval}}$ & $472\,766$\\
interval-certified & $N_{\mathrm{cert}}$ & $252\,282$\\
discarded & $N_{\mathrm{disc}}$ & $485$\\
split & $N_{\mathrm{split}}$ & $219\,999$\\
\textbf{unresolved} & $N_{\mathrm{unres}}$ & $\mathbf 0$\\
float-pass / interval-fail & & $\mathbf 0$\\
$\min_X\bigl(\underline F_{\mathrm{IA}}(X)-T_C\bigr)$ & & $4.628\,931\times10^{-8}$\\
$\max\{\far(X,c_0):X\text{ discarded}\}$ & & $0.001\,978\,039\,980\le\rho_g-10^{-9}$\\
\bottomrule
\end{tabular}
\end{center}
The adaptive segmentation is a disjoint tiling of $\Om$, and $N_{\mathrm{unres}}=0$ means every terminal
box is either certified or discarded. Writing $\mathcal C_{\mathrm{cert}}$ for the union of the certified
boxes and $\mathcal D$ for the union of the discarded ones,
\begin{equation}\label{eq:setdecomp}
\boxed{\;\Om=\mathcal C_{\mathrm{cert}}\;\sqcup\;\mathcal D,\qquad
\mathcal D\subseteq B_1\cup B_2,\qquad F\ge T_C \text{ on }\mathcal C_{\mathrm{cert}}\;}
\end{equation}
Consequently $\Om=\mathcal C_{\mathrm{cert}}\cup B_1\cup B_2$ is a valid cover: the first decomposition of
\eqref{eq:setdecomp} is exact, so no gap is possible, and overlaps among the three pieces are harmless.
\end{proof}

\section{Proof of Theorem~\ref{thm:A}}\label{sec:proof}

Let $x\in\Om$. By Lemma~\ref{lem:global2} and \eqref{eq:setdecomp},
$\Om=\mathcal C_{\mathrm{cert}}\cup B_1\cup B_2$.

\begin{enumerate}
\item If $x\in\mathcal C_{\mathrm{cert}}$, then $F(x)\ge T_C=F(z_*)+10^{-9}>F(z_*)$.
\item If $x\in B_1\cup B_2$, write $x=\sigma^{i}z_*+\delta$ with $i\in\{0,1\}$ and $\|\delta\|\le\rho_g$.
If $\delta=0$, then $F(x)=F(z_*)$ by $F\circ\sigma=F$ (Lemma~\ref{lem:symmetry}). Otherwise
Lemma~\ref{lem:growth2} applied at the corresponding centre gives $F(x)>F(z_*)$.
\end{enumerate}

The two cases cover $\Om$, so $F(x)\ge F(z_*)$ on $\Om$ with equality exactly at $z_*$ and $\sigma z_*$;
hence $C_3=F(z_*)$ and the minimiser set in $\Om$ is exactly $\{z_*,\sigma z_*\}$. The two points are
distinct because $r_2=0.7905\ne0.8302=r_3$. \qed

The hypotheses consumed are: \eqref{eq:deltaA0} from Lemma~\ref{lem:kkt}; the constants of
Lemmas~\ref{lem:sep2}--\ref{lem:growth2}; and the target \eqref{eq:TC} with the partition outcome of
Lemma~\ref{lem:global2}. No property of the Riemann zeta function is used.

\section{The minimiser set in the normalised domain}\label{sec:orbit}

\begin{lemma}[Symmetry]\label{lem:symmetry}
$F\circ\sigma=F$, where $\sigma$ is the paired swap \eqref{eq:swap}.
\end{lemma}

\begin{proof}
Applying $\sigma$ to $x$ merely exchanges the second and third summands of every $S_k$, and the maximum
over $k$ is unchanged.
\end{proof}

\begin{proposition}\label{prop:orbit}
Let $z_*$ be the root of Lemma~\ref{lem:kkt}. Then
\begin{enumerate}
\item[(i)] the raw permutation orbit of $z_*$ under $S_3$ has six points,
$|\mathcal O_{\mathrm{raw}}|=6$;
\item[(ii)] under the normalisation map $\mathcal N$ (re-order the pairs so that the unit-modulus pair
comes first) the orbit collapses to exactly two points,
\[
\bigl|\mathcal N(\mathcal O_{\mathrm{raw}})\bigr|=2,\qquad
\mathcal N(\mathcal O_{\mathrm{raw}})=\{z_*,\sigma z_*\};
\]
\item[(iii)] permuting the three angles while keeping the radii fixed is \emph{not} a symmetry of $F$: for
the certified root,
\[
F(r_2,r_3,\varphi_2,\varphi_1,\varphi_3)=1.126\,876\,598\,534\,850
\;\ne\;0.373\,091\,892\,895\,817=F(z_*) .
\]
\end{enumerate}
\end{proposition}

\begin{proof}
(i) The three pairs $(1,\varphi_1),(r_2,\varphi_2),(r_3,\varphi_3)$ are pairwise distinct because
$r_2\ne r_3$; permuting them gives six distinct raw configurations, and the underlying functional
$\max_k\operatorname{Re}\sum_jz_j^k$ is permutation invariant, so all six carry the same value.
(ii) Of the six permutations, four move the unit-modulus pair into position $2$ or $3$. Re-normalising
(moving the unit pair back to the first position and keeping the relative order of the remaining two)
identifies those four with either $\mathrm{id}$ or $\sigma$ according to whether the relative order is
preserved or reversed. Hence at most two distinct normalised points, and they are distinct since
$r_2\ne r_3$.
(iii) Interchanging $\varphi_1$ and $\varphi_2$ alone merges the undamped term $\cos k\varphi_1$ with the
damped term $r_2^k\cos k\varphi_2$, which is not an invariance; the numerical values differ as stated.
\end{proof}

\section{Independent computational reproduction}

The interval layer of Lemma~\ref{lem:global2} was implemented twice. The second implementation shares only
the (deterministic) floating-point partition and differs structurally: the interval verification is
distributed over three processes with an independently written interval class. It reproduces
\[
N_{\mathrm{cert}}=252\,282,\qquad \text{interval failures}=0,\qquad
\min_X\bigl(\underline F_{\mathrm{IA}}(X)-T_C\bigr)=4.628\,931\,107\,212\,973\,4\times10^{-8},
\]
matching the primary implementation to sixteen significant digits. The two implementations are recorded as
scripts \texttt{dC2iv\_strict.py} (sequential) and \texttt{dC2par\_parallel\_iv.py} (parallel).

\section{Implementation caveats}

\begin{enumerate}
\item \textbf{Facet normals via floating-point SVD.} Lemma~\ref{lem:cover2} obtains the facet normals from
a floating-point SVD. The perturbation allowance in the data ball is
$\varepsilon_{\mathrm{pert}}\approx0.203$, so $10^{-15}\ll\varepsilon_{\mathrm{pert}}$ and the
floating-point error cannot change the sign of $c_X$. This is a numerical-implementation caveat: the lemma
is \emph{not} formalised in interval linear algebra.
\item \textbf{Self-built interval arithmetic.} Lemma~\ref{lem:global2} uses a hand-rolled interval class
over \texttt{mpmath} at $40$ decimal digits, with an outward slack of $10^{-30}$ per operation and $\pi$
enclosed by rational endpoints. An alternative line using \texttt{mpmath.iv} reproduces the bounds
independently.
\item \textbf{\texttt{mpmath} correctness} is assumed, as elsewhere in this project. The interval layer
also reports the float-pass / interval-fail count, which is $0$.
\item \textbf{Lipschitz constants.} The bounds $k\sqrt5$ in Lemma~\ref{lem:sep2} are analytic, not computed.
\item \textbf{$R$} in \eqref{eq:R} is an explicit conservative bound obtained from the $5\times5$ Hessians
with $|\cos|,|\sin|\le1$ and the damping factors retained; since it is an upper bound,
\eqref{eq:locgrowth} is conservative.
\item \textbf{Size of the exclusion radius.} $\rho_g=2c_X/R$ is covering-limited (it is smaller than the
isolation radius $\rho_{\mathrm{iso}}$ of \eqref{eq:rhoiso}); the growth inequality \eqref{eq:locgrowth} is
sharp only in its linear regime, and the margin at $\|\delta\|=\rho_g$ vanishes in the limit.
\end{enumerate}

\section{Boundaries: what is not claimed}

\begin{enumerate}
\item Theorem~\ref{thm:A} concerns $M=3$ in the damped setting. No extrapolation to general $M$, to damped
$M\ge4$, or to the undamped constant $m_M$ is claimed.
\item The value of $C_3$ is known at interval level; the displayed decimal is the upper-bound
configuration's value together with the certified threshold $T_C=F(z_*)+10^{-9}$, not a closed form.
\item No property of the Riemann zeta function is used.
\item Earlier brackets ($0.373\,091\,8\le C_3\le0.373\,092\,075\,762$ and
$0.76\le m_3\le0.764\,081\,100\,745\,854$) are superseded and retained only as checkpoints.
\item The proof is computer-assisted: Lemmas~\ref{lem:kkt}--\ref{lem:global2} are \CA\ statements. The
elementary layer --- Lemma~\ref{lem:symmetry}, Proposition~\ref{prop:orbit}, inequality
\eqref{eq:subadditive}, the Lipschitz bounds and the Taylor step --- is proved by hand in this note.
\end{enumerate}

\appendix

\section{Numerical certificate table}

\begin{center}
\begin{tabular}{lll}
\toprule
item & value & source \\
\midrule
certified root $z_*$ ($\varphi_j/\pi$) & $(0.10911016482862308881,\ 0.8206637095202066754,\ 0.46171937101861024265)$ & \texttt{dB\_krawczyk\_11.py}\\
radii & $r_2=0.79051323395036623846$, $r_3=0.83020729481457293027$ & idem\\
Krawczyk seed / width / $\|I-YJ\|_\infty$ & $10^{-12}$ / $3.93\times10^{-22}$ / $1.9651\times10^{-10}$ & idem\\
$\lambda_{\min}$ ($k=1,2,3,4,5,15$) & $0.111\,860\,12$ & idem\\
$\Delta_{\mathrm{ref}}$, $\rho_{\mathrm{iso}}$ & $0.279\,596\,813\,6$, $4.311\,71\times10^{-3}$ & \texttt{dB5\_local\_growth\_5d.py}\\
$c$, $c_X$, $\rho_{\mathrm{data}}$ & $0.302\,091\,5$, $0.098\,812\,009\,77$, $2.05\times10^{-3}$ & idem\\
$R$, $\rho_g$ & $99.899\,695\,515\,930\,13$, $1.978\,224\,4\times10^{-3}$ & idem\\
$T_C-F(z_*)$ & $10^{-9}$ & \texttt{dC2iv\_strict.py}\\
$N_{\mathrm{eval}},N_{\mathrm{cert}},N_{\mathrm{disc}},N_{\mathrm{split}},N_{\mathrm{unres}}$ & $472\,766$, $252\,282$, $485$, $219\,999$, $0$ & idem\\
$\min_X(\underline F_{\mathrm{IA}}(X)-T_C)$ & $4.628\,931\times10^{-8}$ & idem\\
cross-check & identical to $16$ significant digits & \texttt{dC2par\_parallel\_iv.py}\\
$C_3$ & $0.373\,091\,892\,895\,816\,4\ldots$ & Theorem~\ref{thm:A}\\
\bottomrule
\end{tabular}
\end{center}

\begin{thebibliography}{9}
\bibitem{Andersson} J.~Andersson, \emph{On some power sum problems of Montgomery and Tur\'an},
arXiv:0706.4131; \emph{Lower bounds in some power sum problems}, arXiv:0704.1879.
\bibitem{Montgomery} H.~L.~Montgomery, \emph{Ten Lectures on the Interface between Analytic Number Theory
and Harmonic Analysis}, CBMS Regional Conf.\ Ser.\ Math.\ 84, AMS, 1994 (Chapter 5, Theorem 11).
\bibitem{Palojarvi} L.~Palojarvi, arXiv:1807.01506 (Lemma 2.2 and Theorem 4.1).
\bibitem{Turan} P.~Tur\'an, \emph{On a New Method of Analysis and its Applications},
Wiley--Interscience, 1984.
\end{thebibliography}

\end{document}
